\documentclass[10pt,a4paper]{article}
\usepackage[utf8]{inputenc}
\usepackage[T1]{fontenc}
\usepackage{lmodern}
\IfFileExists{french.ldf}{\usepackage[french]{babel}}{}
\usepackage[margin=19mm,headheight=14pt]{geometry}
\usepackage{amsmath,amssymb,graphicx,booktabs,array,xcolor}
\usepackage{microtype,fancyhdr,tcolorbox}
\AtBeginEnvironment{itemize}{\setlength{\itemsep}{3pt}\setlength{\parskip}{1pt}}
\tcbuselibrary{breakable}
\usepackage[colorlinks=true,linkcolor=blue!45!black,urlcolor=blue!60!black,citecolor=blue!45!black]{hyperref}
\renewcommand{\contentsname}{Sommaire}
\renewcommand{\refname}{Références}
\definecolor{navy}{HTML}{183047}
\definecolor{teal}{HTML}{167D9A}
\definecolor{gold}{HTML}{DBA52D}
\definecolor{pale}{HTML}{F2F6F9}

\pagestyle{fancy}
\fancyhf{}
\fancyhead[L]{\small\textcolor{navy}{Cu\textsubscript{3}Au : ordre et désordre chimique}}
\fancyhead[R]{\small Projet Monte-Carlo}
\fancyfoot[L]{\small Réseau CFC rigide / ensemble canonique}
\fancyfoot[R]{\thepage}
\setlength{\parindent}{0pt}
\setlength{\parskip}{3pt}
\setlength{\emergencystretch}{2em}
\setlength{\textfloatsep}{9pt plus 2pt minus 2pt}
\setlength{\floatsep}{7pt plus 2pt minus 2pt}
\setlength{\intextsep}{8pt plus 2pt minus 2pt}
\setlength{\abovecaptionskip}{4pt}
\setlength{\belowcaptionskip}{2pt}
\setcounter{topnumber}{3}
\setcounter{bottomnumber}{2}
\setcounter{totalnumber}{4}
\renewcommand{\topfraction}{.90}
\renewcommand{\bottomfraction}{.85}
\renewcommand{\textfraction}{.08}
\renewcommand{\floatpagefraction}{.75}
\setcounter{tocdepth}{2}
\newtcolorbox{attention}{breakable,colback=gold!8,colframe=gold!80!black,boxrule=.6pt,
  title={Point d'interprétation},fonttitle=\small}
\newtcolorbox{algocode}[1]{breakable,colback=pale,colframe=navy,boxrule=.6pt,
  title={#1},fonttitle=\bfseries,coltitle=white,
  before={\par\smallskip},after={\par\smallskip}}
\newcommand{\Au}{\mathrm{Au}}
\newcommand{\Cu}{\mathrm{Cu}}
\newcommand{\kb}{k_{\mathrm B}}
\newcommand{\fig}[3]{\begin{figure}[htbp]\centering
\includegraphics[width=#2\linewidth,height=.42\textheight,keepaspectratio]{figures/#1.pdf}
\begingroup\small\caption{#3}\label{fig:#1}\endgroup\end{figure}}
\input{auto_results.tex}
\hypersetup{
  pdftitle={Transition ordre--désordre de Cu3Au},
  pdfauthor={Hugo Vuacheux et Mazarine Lambert},
  pdfsubject={Étude Monte-Carlo sur réseau rigide}
}
\begin{document}
\hypersetup{pageanchor=false}
\begin{titlepage}
\vspace*{5mm}
{\color{teal}\rule{\linewidth}{2pt}}
\vspace{9mm}
{\Large\textcolor{teal}{Modélisation et simulation numérique}\par}
\vspace{8mm}
{\Huge\bfseries\color{navy} Transition ordre--désordre de Cu\textsubscript{3}Au\par}
\vspace{5mm}
{\large Étude Monte-Carlo sur réseau rigide : méthode, résultats et limites\par}
\vspace{5mm}
{\Large\bfseries Hugo Vuacheux et Mazarine Lambert\par}
\vspace{8mm}
\begin{tcolorbox}[colback=pale,colframe=navy,boxrule=.8pt]
\textbf{Système étudié :} $\Lbox^3$ mailles, $N=\Natoms$ sites,
$c_{\Au}=1/4$, interactions de premiers et seconds voisins.\\[2mm]
\textbf{Résultat exact :} $u_0=-3.813\ \mathrm{eV/atome}$.\\[2mm]
\textbf{Maximum de la production canonique L6 :} $T^*_{C_V}=\Tpeak\ \mathrm K$ sur la grille.\\[2mm]
\textbf{Diagnostics de cette production près du maximum :} \DiagnosticStatus.
\IfFileExists{bonus_sizes.tex}{\par\small
\textbf{Complément achevé :} comparaison croisée $L=6,9,12$ en
section~\ref{sec:tailles}. Les diagnostics L9 et L12 près du pic restent
insuffisants ; le maximum canonique L12 observé vaut 1052.5 K.
Les cycles chauffage--refroidissement sont achevés pour $L=\CompletedCycleSizes$.}{}
\end{tcolorbox}
\vfill
\textbf{Traçabilité.} Tous les résultats numériques, tableaux et figures de ce
rapport proviennent des fichiers du dossier de résultats associé. Le code
recalcule les résultats à partir des paramètres du modèle ; aucune valeur
de transition n'est imposée à la simulation.
\vspace{8mm}
\end{titlepage}
\hypersetup{pageanchor=true}
\tableofcontents
\newpage

\section*{Cadre du projet et conventions}
\addcontentsline{toc}{section}{Cadre du projet et conventions}
Dans ce travail, nous étudions le désordre \emph{chimique} d'un alliage Cu$_3$Au. Les sites
du réseau ne se déplacent pas : l'état d'un site est seulement son occupation
par Au ou Cu. La composition et le volume sont fixés. Les températures sont
en kelvins, les énergies en eV et $\kb=8.617333262\times10^{-5}\ \mathrm{eV/K}$.
Dans le code, $s_i=1$ désigne Au et $s_i=0$ Cu. Une paire non orientée est
comptée une seule fois.
Le paramètre $a$ est le côté de la maille conventionnelle. Les longueurs
sont données en unités de $a$ ; une conversion en nanomètres demanderait
une valeur physique de $a$, qui n'est pas imposée dans cette simulation.

\begin{center}
\begin{tabular}{lrrr}
\toprule
Coquille & Distance & Coordination & Énergies de paire (eV)\\
\midrule
Premiers voisins & $a/\sqrt2$ & $12$ & $J_{\Au\Au}=-0.635$\\
 & & & $J_{\Cu\Cu}=-0.582$, $J_{\Au\Cu}=-0.689$\\
Seconds voisins & $a$ & $6$ & $I_{\Au\Au}=I_{\Cu\Cu}=0$\\
 & & & $I_{\Au\Cu}=+0.0063$\\
\bottomrule
\end{tabular}
\end{center}

\subsection*{Repères de vocabulaire}
Un \textbf{site} est une position fixe du réseau ; son occupation est Au ou Cu.
Une \textbf{maille} est le motif géométrique répété dans l'espace. Les quatre
\textbf{sous-réseaux}, repérés par $\mu=0,1,2,3$, distinguent les positions
de base de la maille CFC. Une \textbf{coquille} regroupe les voisins à une
même distance d'un site choisi. Une \textbf{configuration} donne l'espèce
sur chacun des sites à un instant du calcul.

Une \textbf{réplique} est une configuration qui peut voyager entre les
températures. Une \textbf{chaîne} désigne ici une campagne indépendante
avec sa graine et son ensemble de répliques (\Ntemps\ dans la production L6).
Dans le cycle séquentiel, une chaîne est au contraire une seule trajectoire
transmise d'une température à la suivante. Un \textbf{balayage} effectue
$N$ propositions Au--Cu par température, refus compris. L'\textbf{équilibration}
est la phase initiale écartée des moyennes ; la \textbf{production} est la
phase dont on utilise les mesures pour estimer les grandeurs thermiques.
Un \textbf{domaine} est une région spatiale partageant une variante d'ordre ;
ce n'est ni une chaîne ni un sous-réseau. Le symbole $\mu$ est ici un
indice géométrique de sous-réseau, pas un potentiel chimique : aucun terme
de potentiel chimique n'entre dans le test de Metropolis à composition fixée.
Enfin, $L1_2$ est le nom de la structure ordonnée ; $L=12$ est la taille
de la boîte. Les abréviations L6, L9 et L12 du rapport désignent
respectivement $L=6,9,12$, et non trois structures cristallines différentes.

\subsection*{Méthodologie : quatre expériences numériques distinctes}
\addcontentsline{toc}{subsection}{Méthodologie et portée des mesures}
Le Hamiltonien, la composition et les conditions périodiques sont communs,
mais les objectifs et les données ne sont pas interchangeables.
Les résultats de référence des questions 1--9 proviennent de la production
canonique principale L6. La comparaison achevée avec L9 et L12 est regroupée en
section~\ref{sec:tailles}.
Les expériences du bonus sont identifiées séparément ; leurs nombres de
balayages et leurs diagnostics ne remplacent pas ceux de la campagne principale.
\begin{center}\small
\begin{tabular}{>{\raggedright\arraybackslash}p{.23\linewidth}>{\raggedright\arraybackslash}p{.31\linewidth}>{\raggedright\arraybackslash}p{.35\linewidth}}
\toprule
Campagne $L=6$ & Protocole par chaîne et température & Données et usage légitime\\\midrule
Canonique principale & 46 températures, 2 chaînes ; 6000 balayages écartés,
30000 mesurés, pas 10 ; échanges de répliques tous les 5 balayages &
Énergie, ordre, fluctuations et $T^*(L)$ ; états atomiques finaux seulement.\\
Cycle thermique & 46 températures dans chaque sens, 2 trajectoires ;
6000 + 30000 balayages, pas 10 ; aucun échange de répliques &
Mémoire du parcours à budget fini ; un état final par palier pour le GIF.\\
Trempe isotherme & Une trajectoire, consigne $4000\to1000$ K ;
6000 balayages, états tous les 20, aucun transitoire retiré &
Relaxation hors équilibre ; pas de distribution de temps de nucléation.\\
Corrélations canoniques & 14 températures, 2 chaînes ; 3000 + 12000 balayages,
états tous les 20 ; échanges de répliques tous les 5 &
600 états par chaîne et température ; spectres et longueurs conditionnels aux diagnostics.\\
\bottomrule
\end{tabular}\end{center}
Les productions canoniques supplémentaires utilisent le même modèle :
$L=9$ conserve les 46 températures et le budget L6 ; $L=12$ utilise
90 températures, 12000 balayages écartés et 60000 de production par chaîne.
Les captures spatiales de L9 et L12 sont espacées de 100 et 200 balayages.
Ces productions simultanées en température ne sont pas des cycles :
les tailles dont les branches séquentielles sont entièrement achevées
dans cette version sont $L=\CompletedCycleSizes$.

Une graine indépendante crée une nouvelle trajectoire ; les répliques
d'une même chaîne restent couplées par leurs échanges. Les mesures sont
classées par température, après les éventuels échanges. Écarter un nombre
fixé de balayages ne démontre pas l'équilibre : les tests de correction,
la stationnarité observée et l'exploration des variantes répondent à des
questions différentes. Une série gelée peut réussir un test de stabilité
tout en manquant les excitations rares. Les variances thermiques sont
calculées dans chaque chaîne ; une divergence entre chaînes ne doit pas
être convertie en chaleur spécifique.

\textbf{Trois niveaux de certitude.} La géométrie, les différences d'énergie
et les références à $T=0$ sont exactes pour le Hamiltonien donné.
Les moyennes et intervalles Monte-Carlo sont des estimations conditionnelles
à la durée, aux états visités et à la grille. La comparaison à un alliage
réel ajoute une erreur de modèle, qui n'est pas contenue dans les barres
statistiques. Le calcul en micro-eV évite la dérive d'accumulation ; il ne
rend pas les paramètres physiques plus précis que ceux de l'énoncé.

\subsection*{Comment lire les résultats sans mélanger les protocoles}
\begin{itemize}
\item \textbf{Questions 1--9 : référence L6.} Sauf mention contraire,
les figures, les nombres et les diagnostics de ces questions concernent
la production canonique de 864 sites. Les résultats des grandes tailles
ne remplacent pas silencieusement cette référence.
\item \textbf{Bonus : quatre questions complémentaires.} La partie 1
explique le cycle et les GIF L6 ; la partie 2 définit les corrélations
spatiales ; la partie 3 compare les productions canoniques L6/L9/L12 ;
la partie 4 compare leurs cycles séquentiels achevés. Les GIF restent L6.
\item \textbf{Toujours lire valeur, protocole et réserve ensemble.}
$q$ mesure une amplitude globale, $\alpha_1,\alpha_2$ deux voisinages,
$\xi/a$ une portée spectrale. Le temps d'autocorrélation $\tau$ mesure
une dépendance entre mesures successives, en mesures ou en balayages :
ce n'est pas une longueur ni un temps physique en secondes.
\end{itemize}

\section{Question 1 -- Réseau, énergie et algorithme de Metropolis}
\subsection{Construction exacte de la géométrie}
Une maille conventionnelle porte les quatre positions
\[
\mathbf b_0=(0,0,0),\quad
\mathbf b_1=(0,\tfrac12,\tfrac12),\quad
\mathbf b_2=(\tfrac12,0,\tfrac12),\quad
\mathbf b_3=(\tfrac12,\tfrac12,0).
\]
Les sites sont $\mathbf r/a=(i,j,k)+\mathbf b_\mu$, avec
$i,j,k=0,\ldots,L-1$. Ainsi $N=4L^3$, $N_{\Au}=L^3$ et $N_{\Cu}=3L^3$.
La boîte demandée correspond à $L=6$, donc $864$ sites dont $216$ Au.

Le changement de coordonnées $(X,Y,Z)=2\mathbf r/a$ permet de travailler
avec des entiers. Les sites CFC vérifient
\[
X,Y,Z\in\{0,\ldots,2L-1\},\qquad X+Y+Z\text{ pair}.
\]
En effet, les parités possibles sont $(0,0,0)$, $(0,1,1)$, $(1,0,1)$ et
$(1,1,0)$, soit exactement les quatre positions de base. Pour $L=6$, les
coordonnées vont de $0$ à $11$. La coordonnée $12$ est identifiée à $0$.

Les déplacements de premiers voisins sont
\[
(\pm1,\pm1,0),\quad(\pm1,0,\pm1),\quad(0,\pm1,\pm1),
\]
et ceux de seconds voisins sont
\[
(\pm2,0,0),\quad(0,\pm2,0),\quad(0,0,\pm2).
\]
Toutes les coordonnées sont réduites modulo $2L$. Le tableau des voisins
est calculé une fois et reste fixe. Les tailles $L<3$ sont refusées dans cette
implémentation pour éviter les ambiguïtés d'images périodiques multiples.

\fig{01_reseau}{.95}{La maille $L1_2$ et le partage des sites de bord. Les
positions sont celles du réseau CFC, quelle que soit l'occupation chimique.}

\subsection{Énergie et tendance à l'ordre}
\fig{13_coquilles}{.95}{Les deux coquilles autour d'un Au dans la référence
ordonnée $L1_2$. Les douze premiers voisins sont Cu et les six seconds
voisins sont Au. Les coquilles sont définies par leurs distances, même si
leurs occupations changent pendant la simulation.}

L'énergie est
\begin{equation}
E=\sum_{\{i,j\}\in\mathcal P_1}J_{\sigma_i\sigma_j}
+\sum_{\{i,j\}\in\mathcal P_2}I_{\sigma_i\sigma_j}.
\end{equation}
Si l'on utilise les listes de voisins orientées, il faut un facteur $1/2$.
On obtient $Nz_1/2=5184$ paires de premiers voisins et $Nz_2/2=2592$ paires
de seconds voisins pour $N=864$.

Pour comparer les préférences chimiques, on définit
\[
w_1=J_{\Au\Au}+J_{\Cu\Cu}-2J_{\Au\Cu}=0.161\ \mathrm{eV}.
\]
Deux contacts Au--Cu coûtent donc $0.161$ eV de moins qu'un contact Au--Au
et un contact Cu--Cu. Au second voisinage, chaque paire Au--Cu coûte
$0.0063$ eV, tandis qu'une paire identique coûte zéro. L'ordre $L1_2$
satisfait ces deux tendances.

\subsection{Échanges à composition fixée et variation locale}
Un mouvement échange un Au et un Cu, choisis uniformément parmi les paires
d'espèces différentes. Cela conserve exactement la stœchiométrie.
Pour des sites $i$ et $j$, initialement Au et Cu,
\begin{align}
\Delta E={}&\sum_{n=1}^2\sum_{k\in V_n(i)\setminus\{j\}}
[U_n(\Cu,\sigma_k)-U_n(\Au,\sigma_k)]\notag\\
&+\sum_{n=1}^2\sum_{k\in V_n(j)\setminus\{i\}}
[U_n(\Au,\sigma_k)-U_n(\Cu,\sigma_k)],
\end{align}
où $V_n(i)$ est la liste des voisins et $U_1=J$, $U_2=I$ les énergies de paire.
Si les sites sont voisins, leur liaison mutuelle
reste Au--Cu : sa variation est nulle. Elle ne doit pas être recalculée comme
si une seule extrémité changeait.

Le noyau rapide exploite aussi l'identité démontrée à la question 5.
Si $a_i,a_j$ sont les nombres de voisins Au de première coquille, et
$b_i,b_j$ ceux de seconde coquille, en excluant le partenaire échangé,
\begin{equation}
\boxed{\Delta E=0.161(a_j-a_i)+0.0126(b_i-b_j)\quad\text{eV}.}
\end{equation}
Les énergies sont stockées en micro-eV entiers. La somme locale est ainsi
exacte, sans accumulation d'erreurs d'arrondi. Les exponentielles
d'acceptation restent naturellement évaluées en nombres réels.

\subsection{Balance détaillée}
Avec $\beta=1/(\kb T)$, la mesure canonique sur les configurations de composition fixée est
\[
\pi_T(C)=Z^{-1}e^{-E(C)/(\kb T)}.
\]
Comme la proposition d'échange est symétrique, l'acceptation de Metropolis est
\begin{equation}
\boxed{A(C\to C')=\min[1,e^{-\Delta E/(\kb T)}].}
\end{equation}
On a $A(C\to C')/A(C'\to C)=e^{-\beta\Delta E}=\pi(C')/\pi(C)$,
ce qui établit la balance détaillée. Accepter certains mouvements défavorables
est nécessaire pour reproduire les excitations thermiques.

Un balayage représente ici $N$ propositions \emph{Au--Cu}. Les propositions
triviales entre deux espèces identiques sont retirées avant le comptage.
Cette convention modifie l'unité de temps Monte-Carlo, mais pas la mesure
d'équilibre. Les échanges à longue distance sont un outil d'échantillonnage ;
ils ne décrivent pas une cinétique réaliste de diffusion par lacunes.

\fig{14_organigramme}{.72}{Organisation du calcul : proposition, décision,
comptage des tentatives, fin de balayage, échanges puis mesures. Le dessin
résume les étapes ; les règles de proposition, de décision et de mesure sont décrites dans cette section.}

\section{Question 2 -- Ordre à courte portée : paramètres de Cowley}
Pour une référence Au, l'énoncé définit
\[
\alpha_n=1-\frac{p_n^{\Cu\mid\Au}}{c_{\Cu}},\qquad c_{\Cu}=\frac34.
\]
Si $B_n$ est le nombre de paires Au--Cu non orientées de la coquille $n$,
chaque paire possède exactement une extrémité Au. On obtient
\[
p_n^{\Cu\mid\Au}=\frac{B_n}{N_{\Au}z_n},\qquad
\boxed{\alpha_n=1-\frac{B_n}{N_{\Au}z_nc_{\Cu}}.}
\]
Pour la boîte de l'énoncé,
\[
\alpha_1=1-\frac{B_1}{1944},\qquad
\alpha_2=1-\frac{B_2}{972}.
\]
\begin{tcolorbox}[breakable,colback=pale,colframe=teal,title={Du comptage à la probabilité puis à Cowley}]
À première coquille, les $216$ Au donnent $216\times12=2592$ positions
voisines examinées. Si $1944$ de ces positions portent Cu, alors
\[
p_1^{\Cu\mid\Au}=\frac{1944}{2592}=0.75,
\qquad \alpha_1=1-\frac{0.75}{0.75}=0.
\]
La première division mesure une probabilité conditionnelle ; la seconde
la compare à la concentration globale $c_{\Cu}=0.75$.
Si les $2592$ positions portent Cu, $p_1^{\Cu\mid\Au}=1$ et
$\alpha_1=-1/3$. À seconde coquille, on examine $216\times6=1296$
positions : si aucune ne porte Cu, $p_2^{\Cu\mid\Au}=0$ et $\alpha_2=1$.
Le code effectue ce comptage autour de tous les Au, à chaque mesure,
puis moyenne les paramètres obtenus sur la production.
\end{tcolorbox}

Dans $L1_2$, les premiers voisins d'un Au sont tous Cu, donc
$\alpha_1=1-1/(3/4)=-1/3$. Les seconds voisins sont tous Au, donc
$\alpha_2=1$.

Un paramètre négatif indique un excès de voisins d'espèce opposée. Un
paramètre positif indique leur déficit. Zéro est la référence d'un mélange
indépendant pour cette coquille dans la limite d'un grand système.
La nullité de deux paramètres ne prouverait pas, à elle seule, l'absence
de toute corrélation à d'autres distances.

À composition fixée dans une boîte finie, conditionner la référence à être
Au laisse $N_{\Cu}$ Cu parmi $N-1$ autres sites. Ainsi
\[
\left\langle\alpha_n\right\rangle_{T\to\infty}
=1-\frac{N_{\Cu}/(N-1)}{N_{\Cu}/N}
=\boxed{-\frac1{N-1}}.
\]
Pour $N=864$, cette correction vaut environ $-0.00116$.

\fig{12_cowley}{1}{Évolution des paramètres de Cowley $\alpha_1$ et
$\alpha_2$ avec la température. Les barres représentent une erreur-type.
Les lignes horizontales indiquent l'ordre parfait $L1_2$.}

À basse température, les courbes rejoignent les valeurs de l'ordre parfait :
$\alpha_1=-1/3$ et $\alpha_2=1$. Lorsque la température augmente,
les deux paramètres s'éloignent de ces valeurs et se rapprochent de zéro.
Un $\alpha_1$ encore négatif au-dessus de la zone de perte d'ordre global
signifie que les voisins Cu restent favorisés autour d'Au : l'ordre local
peut donc persister après la disparition de l'ordre à longue portée.
La limite aléatoire correspond à $T\to\infty$ ; les mesures à
température finie ne sont pas contraintes à l'atteindre.

\section{Question 3 -- Ordre à longue portée et variantes équivalentes}
Le sous-réseau des sommets comprend $N/4$ sites. En notant
$n_{\Au}^{(0)}$ son nombre d'Au,
\[
p_0=\frac{n_{\Au}^{(0)}}{N/4},\qquad
\boxed{\eta=\frac{p_0-1/4}{3/4}}.
\]
Pour $N=864$, cela s'écrit $\eta=(n_{\Au}^{(0)}-54)/162$.
Dans la variante de référence, $p_0=1$ et $\eta=1$. Dans un mélange
aléatoire, $\langle p_0\rangle=1/4$ et $\langle\eta\rangle=0$.

Il existe quatre variantes équivalentes : Au peut occuper chacun des quatre
sous-réseaux CFC. Dans une autre variante parfaitement ordonnée, $p_0=0$
et $\eta=-1/3$. Pour éviter une mauvaise interprétation, le code conserve
le paramètre du sujet et calcule aussi
\begin{equation}
q=\sqrt{\frac43\sum_{\mu=0}^3(p_\mu-1/4)^2},
\qquad \sum_{\mu=0}^3p_\mu=1.
\end{equation}
On a $q=1$ dans les quatre variantes et $q=0$ pour les occupations moyennes
uniformes. Mais $\langle q\rangle$ est légèrement positif dans une boîte finie
désordonnée, car il mesure une amplitude de fluctuations.
Inversement, un $q$ faible ne prouve pas une occupation indépendante :
plusieurs domaines de variantes différentes peuvent réduire l'ordre global
tout en conservant des voisinages localement ordonnés.

\fig{15_variantes}{1}{Les quatre variantes parfaitement ordonnées sont
équivalentes : seul le sous-réseau qui porte les Au change. $\eta$ dépend
du sous-réseau de référence, tandis que $q=1$ dans les quatre cas. Ces
dessins sont des états de référence construits, pas des mesures Monte-Carlo.}

\begin{center}\small
\begin{tabular}{lccc}\toprule
Configuration & $(p_0,p_1,p_2,p_3)$ & $\eta$ & $q$\\\midrule
Au sur le sous-réseau 0 & $(1,0,0,0)$ & $1$ & $1$\\
Au sur le sous-réseau 1 & $(0,1,0,0)$ & $-1/3$ & $1$\\
Au sur le sous-réseau 2 & $(0,0,1,0)$ & $-1/3$ & $1$\\
Au sur le sous-réseau 3 & $(0,0,0,1)$ & $-1/3$ & $1$\\\bottomrule
\end{tabular}\end{center}

Si ces quatre états sont visités aussi souvent, leur occupation moyenne
vaut $(1/4,1/4,1/4,1/4)$. Pourtant, chaque instantané est ordonné :
\[
\boxed{\langle q(p)\rangle=1\quad\text{alors que}\quad
q(\langle p\rangle)=0.}
\]
Les courbes de gauche de la figure suivante sont des moyennes sur les
mesures \emph{et les chaînes}. Leur proximité avec $1/4$ à basse température
reflète la visite de plusieurs variantes ; elle ne suffit pas à conclure
au désordre. Il faut les lire avec $\langle q\rangle$, qui reste proche de 1.

\fig{09_sous_reseaux}{1}{Paramètre signé de l'énoncé, amplitude invariante et
occupations des sous-réseaux. Une moyenne entre variantes ne doit pas être
confondue avec une configuration instantanée désordonnée.}

\section{Question 4 -- Vérifications du programme}
Le dossier contient une suite de tests automatiques : \TestCount\ tests
ont été relancés avec succès pour la présente mise à jour. Leur date,
leurs valeurs de référence et leur environnement sont précisés ci-dessous.
Ils
contrôlent la géométrie périodique, la réciprocité des voisinages, les quatre
variantes, les valeurs analytiques, les différences d'énergie locales et
globales, la conservation de composition et la reproductibilité.
Les tests supplémentaires vérifient le bootstrap du maximum, sa
reproductibilité et son invariance à une constante ajoutée à l'énergie.
Ils vérifient également la continuité du cycle thermique et l'identité
des trajectoires après reprise à un palier sauvegardé.
Les tests de checkpoints vérifient aussi la reprise exacte au milieu
d'une campagne canonique et le refus d'une configuration incompatible.
Trois tests de l'analyse croisée contrôlent l'identité énergie--Cowley,
la décomposition avec covariance sans ajout de dispersion entre chaînes,
et le comptage des populations sans franchir artificiellement les
frontières entre chaînes.
\input{current_validation.tex}

\subsection{Création d'une paire d'antisites}
Échangeons un Au et un Cu premiers voisins dans l'état $L1_2$. En excluant
leur liaison mutuelle, le premier site a onze voisins Cu, et le second trois
Au et huit Cu. La variation de première coquille vaut
\begin{align*}
\Delta E_1={}&11(J_{\Cu\Cu}-J_{\Au\Cu})
+3(J_{\Au\Au}-J_{\Au\Cu})\\
&+8(J_{\Au\Cu}-J_{\Cu\Cu})=3w_1=0.483\ \mathrm{eV}.
\end{align*}
Douze liaisons identiques de seconde coquille deviennent Au--Cu :
\[
\Delta E_2=12(0.0063)=0.0756\ \mathrm{eV}.
\]
Ainsi
\[
\boxed{\Delta E_{\text{échange 1er voisin}}=0.5586\ \mathrm{eV}.}
\]
L'échange inverse rend exactement l'énergie opposée.

\subsection{Contrôles indépendants}
Le noyau rapide par comptage est confronté à une somme explicite des
matrices de paire pour des échanges aléatoires, voisins et non voisins, sur
plusieurs tailles. Un petit ensemble restreint de configurations est énuméré
exactement pour vérifier la stationnarité canonique de la matrice de
Metropolis. Le signe de l'acceptation des échanges de répliques est également
vérifié par balance détaillée.

Pendant la campagne, à la fin de chaque segment, le code recalcule l'énergie
complète de chaque réplique et impose l'égalité entière avec l'énergie suivie
localement. Il vérifie également la composition et la borne $E\ge E_0$.
Une incohérence arrête immédiatement le calcul.

\section{Question 5 -- Énergie de cohésion à température nulle}
Chaque Au a douze voisins Cu. Il y a donc $216\times12=2592$ paires Au--Cu,
aucune paire Au--Au de première coquille, et $5184-2592=2592$ paires Cu--Cu.
Les seconds voisins sont de même espèce et ne contribuent pas à l'énergie.
Le tableau compte les paires \emph{non orientées}, une seule fois chacune :
\begin{center}\small
\begin{tabular}{lrrrr}\toprule
Coquille & Au--Au & Au--Cu & Cu--Cu & Total\\\midrule
Première & 0 & 2592 & 2592 & 5184\\
Deuxième & 648 & 0 & 1944 & 2592\\\bottomrule
\end{tabular}\end{center}
Par exemple, les $216\times6$ contacts Au--Au de seconde coquille sont
comptés depuis les deux extrémités, soit $216\times6/2=648$ paires.
\[
E_0=2592(-0.689)+2592(-0.582)=-3294.432\ \mathrm{eV}.
\]
Donc
\begin{equation}
\begin{aligned}
u_0&=\frac{E_0}{N}=\frac{-3294.432}{864}
=\boxed{-3.813\ \mathrm{eV/atome}},\\
E_0^{\Cu_3\Au}&=\frac{E_0}{N/4}=\frac{-3294.432}{216}=4u_0
=\boxed{-15.252\ \mathrm{eV/formule}}.
\end{aligned}
\end{equation}
Une unité de formule Cu$_3$Au contient trois Cu et un Au, donc quatre
atomes ; la boîte contient $864/4=216$ unités de formule.
La cohésion exprimée comme énergie positive de dissociation vaut
$3.813\ \mathrm{eV/atome}$, relativement au zéro énergétique du modèle.
Cette cohésion se distingue de l'\textbf{énergie de formation}, référencée
aux cristaux de Cu pur et d'Au pur. Dans le même modèle CFC,
$u_{\Cu}^{\rm pur}=6J_{\Cu\Cu}=-3.492$ et
$u_{\Au}^{\rm pur}=6J_{\Au\Au}=-3.810$ eV/atome. Ainsi
$\Delta u_f=u_0-[3u_{\Cu}^{\rm pur}/4+u_{\Au}^{\rm pur}/4]
=-0.2415$ eV/atome, soit $-0.966$ eV par formule Cu$_3$Au.
Ce signe favorise l'alliage ordonné par rapport à ces deux cristaux
séparés ; il ne détermine pas à lui seul sa stabilité à température finie.

\subsection{Preuve que \texorpdfstring{$L1_2$}{L12} est un état fondamental}
Soient $A_1,B_1,C_1$ les nombres de paires Au--Au, Au--Cu et Cu--Cu de
première coquille et $M_1=Nz_1/2$ leur nombre total. À composition fixée,
\[
2A_1+B_1=12N_{\Au},\qquad A_1+B_1+C_1=M_1.
\]
En éliminant $B_1$ et $C_1$ dans l'énergie, on trouve
\[
E_1=J_{\Cu\Cu}M_1+(J_{\Au\Cu}-J_{\Cu\Cu})12N_{\Au}+w_1A_1.
\]
Pour les seconds voisins, $E_2=0.0063B_2$. Puisque $L1_2$ possède
$A_1=B_2=0$,
\begin{equation}
\boxed{E-E_0=0.161A_1+0.0063B_2\ge0.}
\end{equation}
C'est une preuve du minimum global pour les paramètres donnés. Elle ne
nécessite aucune simulation.

\section{Question 6 -- Équilibrer à différentes températures}
La simulation doit échantillonner
\[
P_T(C)=\frac{e^{-\beta E(C)}}{\sum_{C'\,\text{de même composition}}e^{-\beta E(C')}}.
\]
Une configuration finale seule n'est pas une moyenne thermique. À chaque
température, on distingue une phase d'équilibration, écartée de l'analyse,
et une phase de production.

\subsection{Échanges de répliques et mesure canonique}
Pour faciliter les passages entre états métastables, plusieurs répliques
évoluent à des températures $T_r$. La mesure jointe est le produit
$\Pi(\{C_r\})\propto\prod_r e^{-\beta_rE(C_r)}$. L'échange de deux
configurations à températures voisines est accepté avec
\begin{equation}
A_{r,r+1}=\min\left[1,
\exp\{(\beta_r-\beta_{r+1})(E_r-E_{r+1})\}\right].
\end{equation}
Le rapport des poids après et avant échange donne directement cet exposant.
La marginale à chaque température reste canonique. La mesure doit être
classée par \emph{température fixe}, et non par identité de réplique.

Une session d'échange a lieu tous les \ExchangeEvery\ balayages. Le code
alterne les liens pairs et impairs : un lien donné est donc testé toutes
les deux sessions, soit tous les 10 balayages dans cette campagne. Les noyaux de
Metropolis restent ceux demandés par l'énoncé ; les échanges de répliques
ajoutent un mécanisme d'échantillonnage qui conserve la mesure cible
\cite{metropolis,hukushima}.

\subsection{Protocole réellement exécuté}
\begin{center}
\begin{tabular}{lr}
\toprule
Paramètre & Valeur\\
\midrule
Nombre de sites & \Natoms\\
Nombre de températures & \Ntemps\\
Chaînes indépendantes & \Nchains\\
Balayages d'équilibration par chaîne & \BurnSweeps\\
Balayages de production par chaîne & \ProductionSweeps\\
Intervalle entre mesures & \SampleEvery\\
Mesures par température et par chaîne & \Ndraws\\
Maximum sur la grille & $\Tpeak\ \mathrm K$\\
\bottomrule
\end{tabular}
\end{center}
Les graines et toutes les températures figurent dans le fichier de
configuration et dans \texttt{summary.json}. Les chaînes démarrent dans
des variantes ordonnées différentes ; dans la seconde chaîne, les répliques
chaudes démarrent désordonnées. Un départ entièrement aléatoire à très basse
température risquerait de se figer avant d'atteindre l'état fondamental.

\fig{16_coupes_2D}{.95}{Deux plans adjacents des mêmes configurations de la
première chaîne. Les deux hauteurs sont nécessaires pour montrer les quatre
sous-réseaux : une seule coupe peut ne contenir aucun Au dans une variante
pourtant parfaitement ordonnée. L'ordre est plus lisible ici que dans le
nuage tridimensionnel, conservé comme fichier complémentaire
\texttt{07\_configurations.pdf} sans seconde reproduction des mêmes états.}
\fig{17_moyennes_cumulees}{.95}{Comparaison des moyennes cumulées de
production dans les deux chaînes à $\Tpeak$ K. Une moyenne cumulée utilise
toutes les mesures précédentes ; ses points ne sont donc pas indépendants.}

Pour $M'$ mesures de production, la courbe cumulée est
$\overline A(M')=(1/M')\sum_{m=1}^{M'}A_m$. Sa stabilisation et l'accord
des chaînes complètent les traces instantanées, sans prouver à eux seuls
l'absence de bassins métastables non visités. Les contrôles de circulation
donnent 4 puis 3 allers-retours complets froid--chaud--froid pour les deux
chaînes de cette campagne ; leur faible nombre reste une limite à considérer.
\fig{10_repliques}{1}{Trajectoires de quelques identités de répliques.
La circulation en température complète les contrôles locaux d'équilibre.}

\section{Question 7 -- Grandeurs physiques et incertitudes}
Les moyennes numériques approchent
$\langle A\rangle_T=\sum_C A(C)P_T(C)$. On calcule l'énergie par atome,
les deux paramètres de Cowley, $\eta$, $q$ et les occupations des quatre
sous-réseaux.

\subsection{Capacité calorifique et fluctuations d'ordre}
\input{representative_results.tex}
Ce tableau est calculé à partir du CSV des mesures. La valeur observée
$c_V=0$ au froid signifie qu'aucune fluctuation d'énergie n'a été mesurée
dans cette série ; elle ne démontre pas que la capacité calorifique exacte
est strictement nulle à une température finie.

La capacité calorifique mesure la variation d'énergie moyenne quand on
chauffe à volume constant :
\[
C_V=\left(\frac{\partial\langle E\rangle}{\partial T}\right)_V,
\qquad c_V=C_V/N.
\]
Le déroulement concret est
\[
\text{configurations Metropolis}\ \longrightarrow\
\{E_1,\ldots,E_M\}\ \longrightarrow\
\langle E^2\rangle-\langle E\rangle^2\ \longrightarrow\ c_V.
\]
Metropolis fournit les configurations avec leurs poids thermiques ; le
post-traitement utilise les énergies mesurées, pas le taux d'acceptation.
Dans l'ensemble canonique,
$Z=\sum_C e^{-\beta E(C)}$ et $\langle E\rangle=-\partial_\beta\ln Z$.
En dérivant une seconde fois, on obtient la variance d'énergie.
À Hamiltonien indépendant de la température,
\[
\frac{\partial\langle E\rangle}{\partial\beta}
=-\big(\langle E^2\rangle-\langle E\rangle^2\big).
\]
Comme $d\beta/dT=-1/(\kb T^2)$,
\begin{equation}
\boxed{\frac{c_V}{\kb}=
\frac{\langle E^2\rangle-\langle E\rangle^2}{N(\kb T)^2}.}
\end{equation}
Il faut utiliser l'énergie \emph{totale} dans cette formule. La contribution
calculée est configuratoire : elle n'inclut pas la chaleur des vibrations.

On suit aussi
\[
X_q=\frac{N}{\kb T}\big(\langle q^2\rangle-\langle q\rangle^2\big),
\quad
U_q=1-\frac{\langle q^4\rangle}{3\langle q^2\rangle^2}.
\]
$X_q$ est une réponse à un champ artificiel couplé comme $-Nhq$ ; il ne
faut pas l'identifier sans précaution à une susceptibilité d'un composant
du paramètre d'ordre. Le cumulant est fourni comme diagnostic complémentaire.
Le facteur $3$ est une convention scalaire ; sa limite désordonnée n'est
pas supposée nulle pour cette amplitude à plusieurs composantes.
Par exemple, si les trois composantes de l'ordre sont gaussiennes,
indépendantes et de même variance, leurs moments donnent
$\langle q^4\rangle/\langle q^2\rangle^2=5/3$, donc $U_q=4/9$.
Cette référence explique la valeur chaude proche de 0.44, sans supposer
que toutes les distributions à température finie soient gaussiennes.
$X_q$ est exprimé en eV$^{-1}$ ; $q$, $U_q$ et $c_V/\kb$ sont sans dimension.

\subsection{Autocorrélation, erreurs et comparaison des chaînes}
Les mesures successives ne sont pas indépendantes. Pour une série $A_t$,
\[
\tau_{\mathrm{int}}=\frac12+\sum_{\ell\ge1}\rho_A(\ell),\qquad
N_{\mathrm{eff}}\simeq\frac{M}{2\tau_{\mathrm{int}}}.
\]
L'\textbf{ESS} est ce nombre estimé d'informations indépendantes ;
ce n'est pas le nombre de lignes enregistrées. Le diagnostic
$\widehat R$ compare les distributions des chaînes et de leurs moitiés :
une valeur proche de 1 indique un accord observé, tandis qu'une valeur
élevée signale des moyennes ou des dispersions différentes.
Le \textbf{bootstrap par blocs} rééchantillonne des morceaux contigus
de trajectoire afin de conserver une partie de leur dépendance temporelle.
Une erreur-type, notée SE dans les tableaux, mesure l'incertitude estimée
de la grandeur calculée ; ce n'est pas la fluctuation thermique de cette grandeur.
Le code utilise une troncature par paires positives de l'autocorrélation
et un bootstrap par blocs circulaires. Les blocs ont une longueur d'au moins
$2\tau_{\mathrm{int}}$, estimée à partir de l'énergie et de $q$, sauf lorsque
la série est trop courte pour le permettre ; ce cas est signalé par le faible
nombre de blocs. Les erreurs retiennent aussi la dispersion entre chaînes.
Les variances thermiques sont calculées \emph{dans chaque chaîne} avant
moyennage pour éviter de transformer une divergence entre chaînes en
capacité calorifique artificielle.

Les contrôles indicatifs demandent $\widehat R<1.05$ pour l'énergie et $q$,
un effectif indépendant estimé d'au moins $200$, au moins $20$ blocs par
chaîne, et un écart entre demi-séries inférieur à trois unités d'erreur
estimée. Les graines indépendantes ne rendent pas les températures d'une
même campagne indépendantes.
Ces seuils sont des filtres d'alerte. Deux chaînes restées dans la même
population métastable peuvent présenter un bon accord et manquer ensemble
une autre population. Réussir les contrôles locaux ne certifie donc pas
les poids de toutes les configurations à l'équilibre.

Lorsque E et q sont constants dans une série froide, le code attribue
$\tau=1/2$ et peut obtenir $\widehat R=1$ : ce traitement numérique ne
prouve pas que toutes les excitations ou variantes ont été explorées.
Le choix des blocs utilise E et q ; il reste utile de tester également
la corrélation des paramètres de Cowley et la stabilité des erreurs
lorsque la longueur des blocs est augmentée.

Les paramètres de Cowley sont représentés une seule fois en
figure~\ref{fig:12_cowley}. L'énergie, $q$ et les réponses au voisinage de la
transition sont regroupés en figure~\ref{fig:11_zoom_transition} ; les
occupations et $q$ sur la gamme complète figurent en
figure~\ref{fig:09_sous_reseaux}. Les barres donnent une erreur-type
estimée, pas une enveloppe de l'erreur de modèle. Les exports
\texttt{02\_ordres.pdf} et \texttt{03\_thermodynamique.pdf} restent disponibles,
sans reproduire une deuxième fois ces courbes dans le PDF.
\fig{05_diagnostics}{1}{Diagnostics par température. Un indicateur favorable
ne prouve pas l'exploration de tous les états métastables.}
\fig{06_traces}{1}{Énergie et ordre à température fixe. La phase initiale
grisée est exclue des statistiques.}

\subsection{Vérification de la limite de haute température}
Dans un mélange indépendant de grande taille,
$P(\Au\Au)=1/16$, $P(\Cu\Cu)=9/16$ et $P(\Au\Cu)=3/8$. Ainsi
\begin{align*}
u_\infty={}&6\left(\frac{J_{\Au\Au}}{16}
+\frac{9J_{\Cu\Cu}}{16}+\frac{3J_{\Au\Cu}}8\right)
+3\left(\frac{3I_{\Au\Cu}}8\right)\\
={}&-3.7455375\ \mathrm{eV/atome}.
\end{align*}
À $N$ fini, les probabilités exactes sont
\[
P(\Au\Au)=\frac{N_{\Au}(N_{\Au}-1)}{N(N-1)},\quad
P(\Cu\Cu)=\frac{N_{\Cu}(N_{\Cu}-1)}{N(N-1)},\quad
P(\Au\Cu)=\frac{2N_{\Au}N_{\Cu}}{N(N-1)}.
\]
Le code utilise ces corrections pour les tests. La limite exacte pour la
boîte étudiée vaut $\RandomLimit\ \mathrm{eV/atome}$. Même à plusieurs
milliers de kelvins, les interactions peuvent encore créer des corrélations :
il ne faut pas forcer les courbes à rejoindre la limite $T\to\infty$.

\section{Question 8 -- Repérer la transition et comparer au champ moyen}
\subsection{Résultat Monte-Carlo et résolution de la grille}
Le maximum de capacité calorifique sur les températures échantillonnées est
\begin{equation}
\boxed{T^*_{C_V}(L=\Lbox)=\Tpeak\ \mathrm K.}
\end{equation}
Les deux températures voisines de ce maximum sont $\Tleft$ et $\Tright$ K.
L'intervalle $[\Tleft,\Tright]$ K décrit le voisinage sur la grille :
\textbf{ce n'est pas un intervalle de confiance statistique}. Il ne couvre
ni une éventuelle erreur systématique d'équilibration ni le déplacement
avec la taille.

Les maxima trouvés séparément dans les chaînes sont : \ChainPeaks.
Les diagnostics des cinq températures centrées sur le maximum sont résumés
ci-dessous. Leur statut global est : \DiagnosticStatus.
\input{critical_table.tex}
\fig{11_zoom_transition}{1}{Agrandissement de la zone de transition :
ordre, énergie et fluctuations. La zone grisée représente les voisins de
grille, sans interprétation comme intervalle de confiance.}

Une boîte finie n'a pas de singularité thermodynamique. Une seule taille
permet de repérer un maximum apparent, mais pas d'extrapoler $T_c$ à la limite
$N\to\infty$. La comparaison des trois tailles en
section~\ref{sec:tailles} examine distributions et déplacements des maxima,
avec les réserves d'échantillonnage nécessaires \cite{binder}.

\fig{04_distributions}{1}{Distributions près de la zone de transition.
L'inspection de la coexistence complète celle du maximum de capacité
calorifique ; les histogrammes restent sensibles à la durée du calcul.}

\subsection{Incertitude statistique du maximum et concordance des indices}
La hauteur du pic possède une erreur-type, mais celle-ci n'est pas une
erreur en kelvins sur sa position. On estime cette dernière par un
bootstrap \emph{conjoint} des températures : dans chaque chaîne, les mêmes
indices de blocs sont appliqués simultanément à toutes les températures.
Les deux chaînes sont rééchantillonnées séparément. Pour chaque tirage,
on recalcule les variances d'énergie dans chaque chaîne, on moyenne leurs
capacités calorifiques puis on recherche le maximum sur la grille initiale.
Cette synchronisation conserve les dépendances entre températures induites
par les échanges de répliques.

La longueur de référence vaut
$b=\lceil2\max_{T,c,A\in\{E,q\}}\tau_{c,A}\rceil$, avec les bornes
numériques décrites dans la partie statistique (plafond de $M/4$ mesures). On la compare à une longueur environ deux
fois plus petite et deux fois plus grande. Les résultats sont obtenus
uniquement à partir des archives de production :
\input{peak_uncertainty.tex}
\fig{18_incertitude_pic}{.95}{Fréquence de la position du maximum parmi
les tirages bootstrap pour trois longueurs de bloc. Les températures
possibles sont celles de la grille : aucune précision sous la maille de
température n'est créée par ce calcul.}

Un écart-type inférieur au pas de 5 K traduit ici des changements rares
de la case de maximum ; il ne localise pas le pic entre les températures
simulées. L'intervalle percentile est une estimation statistique conditionnelle à
la taille, aux deux chaînes observées et à la grille. Il ne couvre ni
le déplacement avec $L$, ni les états métastables éventuellement non
visités. S'il se réduit à un point, cela ne signifie pas que la vraie
température de transition est connue exactement. Les changements de pente
et les maxima de fluctuations ne définissent pas exactement le même
critère : leur proximité identifie une zone cohérente, pas plusieurs
mesures indépendantes d'une température critique unique.

Une étude de plusieurs tailles à précision statistique contrôlée
renforcerait l'interprétation thermodynamique. Les résultats de cette
question concernent $L=6$.
Les campagnes $L=9$ et $L=12$ sont achevées et techniquement validées ;
leurs résultats, avec leurs réserves statistiques, sont présentés en
section~\ref{sec:tailles}.

\subsection{Estimation analytique de Bragg-Williams}
Pour une famille d'états à un sous-réseau privilégié,
avec $0\leq\eta\leq1$ dans cette comparaison d'amplitudes,
\[
p_0=\frac{1+3\eta}{4},\qquad
p_1=p_2=p_3=p_f=\frac{1-\eta}{4}.
\]
Le champ moyen néglige les corrélations entre occupations. Les premiers
voisins relient des sous-réseaux distincts ; les seconds voisins appartiennent
au même sous-réseau. Le comptage donne
\[
\frac{B_1}{N}=\frac94+\frac34\eta^2,\qquad
\frac{B_2}{N}=\frac98(1-\eta^2).
\]
L'énergie moyenne s'écrit
\[
u_{\mathrm{MF}}(\eta)=u_\infty-A\eta^2,
\quad A=\frac38w_1-\frac9{16}w_2=0.0674625\ \mathrm{eV},
\]
avec $w_2=I_{\Au\Au}+I_{\Cu\Cu}-2I_{\Au\Cu}=-0.0126$ eV.

L'entropie de mélange par site d'un sous-réseau vaut $-\kb h(p)$, où
$h(p)=p\ln p+(1-p)\ln(1-p)$ et $0\ln0=0$. Donc
\[
f_{\mathrm{MF}}(\eta,T)=u_\infty-A\eta^2+\kb T g(\eta),
\quad g(\eta)=\tfrac14h(p_0)+\tfrac34h(p_f).
\]
La coexistence vérifie à la fois la stationnarité du minimum ordonné et
l'égalité de son énergie libre à celle du minimum désordonné :
\[
2A\eta_t=\kb T_tg'(\eta_t),\qquad
A\eta_t^2=\kb T_t[g(\eta_t)-g(0)].
\]
En éliminant $T_t$,
\[
\eta_tg'(\eta_t)=2[g(\eta_t)-g(0)],\qquad
g'(\eta)=\frac3{16}\ln\frac{(1+3\eta)(3+\eta)}{3(1-\eta)^2}.
\]
La résolution donne $\eta_t\simeq0.46358$ et
\[
\boxed{T_t^{\mathrm{MF}}\simeq1712.18\ \mathrm K.}
\]
L'annulation de $f''(0,T)$ donne seulement la perte de stabilité locale de
l'état désordonné : $T_{\mathrm{sp}}=2A/\kb\simeq1565.74$ K. Cette
spinodale est distincte de la coexistence. Le saut d'ordre prédit ici est
une propriété de l'approximation de champ moyen, et ne remplace pas une
étude de la nature de la transition du modèle spatial.

\fig{08_champ_moyen}{1}{Énergies libres approchées et comparaison entre
amplitudes d'ordre. L'estimation de champ moyen n'est pas utilisée pour
fixer le résultat Monte-Carlo.}

\section{Question 9 -- Discussion des résultats et limites}
\input{discussion_results.tex}

Le mécanisme général est une compétition entre énergie et entropie :
$F=U-TS$. L'ordre minimise l'énergie, tandis que le nombre très élevé de
configurations désordonnées favorise leur poids statistique à haute
température. La disparition de l'ordre global peut précéder la disparition
des préférences chimiques locales : un $q$ faible ne contraint pas
$\alpha_1$ et $\alpha_2$ à être exactement nuls.

Les limites principales sont explicites : tailles finies, durées de
simulation finies, grille de températures discrète, et approximation physique
du réseau rigide. Les vibrations, la dilatation et les relaxations autour des
antisites sont absentes. Les constantes de paire sont supposées indépendantes
de la température. On ne présente donc pas la température du modèle comme
une prédiction expérimentale exacte.

\subsection{Comparaison critique avec \texorpdfstring{Cu$_3$Au}{Cu3Au} réel}
Les mesures de calorimétrie et de diffusion situent le désordre chimique de
Cu$_3$Au près de 660--663 K, et non à 1050 K
\cite{benisek,paras}. L'écart d'environ 390 K est de l'ordre de 60\,\% :
le maximum du présent modèle est donc qualitativement pertinent mais
quantitativement non calibré au matériau. Passer de $L=6$ à $L=9$ ou $12$ ne
constitue pas une correction des interactions du Hamiltonien. On ne peut
pas attribuer cet écart à la seule taille finie, ni aux seules vibrations,
sans calculs de contrôle séparant ces contributions.
Un maximum calorimétrique expérimental dépend aussi du protocole :
Benisek et Dachs rapportent un pic à 675 K pour leur échantillon le
plus ordonné et une réaction vers 660 K pour des échantillons initialement
désordonnés \cite{benisek}. Il faut donc distinguer température de
transition de référence et maximum mesuré pendant une histoire thermique.
La transition réelle est décrite comme du premier ordre
\cite{paras}. L'arrondi des courbes dans notre boîte finie ne démontre
donc pas une transition continue, et l'on ne doit pas exiger une divergence
critique de $\xi$ pour conclure à l'existence d'une transition.

\begin{center}\small
\begin{tabular}{>{\raggedright\arraybackslash}p{.22\linewidth}>{\raggedright\arraybackslash}p{.34\linewidth}>{\raggedright\arraybackslash}p{.33\linewidth}}
\toprule
Observation ou attente physique & Ce que donne le modèle & Limite ou interprétation\\\midrule
Ordre chimique au froid & $L1_2$, quatre variantes, défauts d'antisites et
préférence pour les liaisons Au--Cu & Accord structurel et énergétique
pour ce modèle ; les positions atomiques ne se relaxent pas.\\
Ordre local après perte d'ordre global & $q$ devient faible alors que
$\alpha_1,\alpha_2$ restent mesurables & Accord qualitatif avec la diffusion
expérimentale sensible à l'ordre à courte portée \cite{owen}.\\
Chaleur mesurée en calorimétrie & Pic de la seule contribution
configuratoire $c_V$ & Une expérience mesure généralement $C_P$, avec
vibrations et autres contributions ; ni le fond thermique, ni le pic
total ne sont prédits ici \cite{benisek}.\\
Diffusion et croissance de domaines & Mise en ordre rapide avec des
échanges Au--Cu à toute distance & Pas de lacunes ni de barrières de saut
calibrées ; aucun temps physique ni coefficient de diffusion ne peut
être déduit du GIF \cite{kessler}.\\
Volume, surfaces et interfaces & Boîte périodique, volume fixé et interfaces
chimiques possibles & Pas de surface libre, joint de grain ou contrainte
élastique relaxable ; les images périodiques peuvent interagir.\\
Très haute température & Le réseau reste artificiellement solide même
à 4000 K & Cette borne sert à explorer la limite mathématique désordonnée ;
le modèle n'a aucun mécanisme de fusion et ne représente pas un solide réel à cette température.\\
\bottomrule
\end{tabular}\end{center}

\textbf{Ce qui est observé malgré les simplifications.} La compétition
énergie--entropie, les quatre variantes, la création de défauts, la perte
d'ordre global et la persistance de préférences locales sont des mécanismes
capturés par les occupations et les interactions de paires. Leur présence
ne valide pas automatiquement leurs amplitudes, leur température ou
leur vitesse dans le matériau réel.

\textbf{Ce que les animations ne résolvent pas.} Selon la température de
trempe et la dynamique choisie, un modèle d'alliage peut présenter une
incubation par nucléation ou une mise en ordre plus homogène
\cite{kessler}. Une seule trajectoire, sans temps d'attente apparent à
la résolution de 20 balayages, ne permet pas de choisir entre ces scénarios.
Les échanges globaux peuvent faciliter une réorganisation collective que
la diffusion locale réaliserait autrement. Inversement, un pic ou une
bosse d'occupation de sous-réseau peut provenir d'un changement de variante,
sans constituer un germe géométrique nouveau.

\subsection{Approximations statistiques et spatiales à ne pas confondre}
Un $q$ résiduel au chaud n'est pas nécessairement un défaut : $q$ est une
norme positive, dont la moyenne ne s'annule pas exactement dans une petite
boîte désordonnée. Une absence de fluctuations au froid borne seulement
ce qui a été vu pendant cette production. Les erreurs par blocs corrigent
la dépendance temporelle observée ; elles ne reconstruisent pas un bassin
métastable jamais visité.

Les longueurs du bonus sont celles d'une enveloppe d'ordre définie par
maille, avec moyenne sur les trois composantes et sur les directions
équivalentes. Ce choix fournit un indicateur scalaire ; il peut masquer
une anisotropie des domaines. Le premier vecteur accessible à $L=6$ vaut
$k_{\min}a\simeq1.047$, loin d'une résolution continue de très petits $k$.
Les distances périodiques et la petite gamme de $r$ interdisent un ajustement
robuste d'une décroissance exponentielle ou d'un exposant critique.
Une longueur canonique supérieure à la boîte peut refléter le mélange
des variantes et l'ordre uniforme ; ce n'est pas la taille d'un domaine
mesuré dans les coupes. La longueur intra-variante est un diagnostic
conditionnel, sans intervalle de confiance calculé dans cette étude.

\subsection{Ce que les tailles supplémentaires permettent de tester}
Le côté passe de $6a$ à $9a$, soit un facteur 1.5 ; le nombre de sites
passe de 864 à 2916, soit un facteur 3.375, avec 729 Au et 2187 Cu.
Le plus petit décalage réciproque passe de $1.047/a$ à $0.698/a$.
Cette augmentation peut réduire l'arrondi des pics et l'influence des
images périodiques, et offrir plus de place aux domaines. Elle ne garantit
aucune amélioration de l'équilibre : les temps de relaxation et les
barrières peuvent augmenter, tandis que les échanges de répliques
deviennent moins efficaces à grille inchangée.

La comparaison L6/L9 conserve le Hamiltonien, les 46 températures,
6000 balayages initiaux et 30000 de production par chaîne. Les moyennes
sont comparées par atome ; les maxima $T^*(L)$, leurs intervalles,
les histogrammes, R-hat, ESS et la circulation des répliques sont
examinés ensemble en section~\ref{sec:tailles}. Les valeurs de $q$ ne sont pas comparées comme
si leur biais de norme finie était indépendant de $N$. Des états
atomiques espacés de 100 balayages ont été conservés pour permettre
une analyse spatiale après le calcul. Ce budget identique de balayages
n'est pas une garantie de nombres identiques d'états indépendants.

L12 porte ensuite la boîte à 6912 sites, avec un plus petit vecteur
accessible $k_{\min}a=0.524$. Sa grille est densifiée et sa production
doublée ; les effets de taille et de protocole ne sont donc pas séparés
parfaitement. Les trois tailles permettent d'observer des tendances,
mais leurs diagnostics critiques inégaux ne justifient ni extrapolation
quantitative ni loi universelle. L'analyse croisée de la
section~\ref{sec:tailles} détaille ce que l'augmentation améliore et
les défauts d'échantillonnage qu'elle n'a pas supprimés.
L'absence de phonons, de dilatation, de lacunes et de dynamique de diffusion
reste la même pour les trois tailles.

Une distribution à deux pics est compatible avec une coexistence ; une
hystérésis est compatible avec une métastabilité. Aucune de ces observations
isolées n'est une preuve suffisante de l'ordre thermodynamique de la
transition. Il faudrait renforcer les statistiques et comparer plusieurs
tailles pour conclure sur ce point.

\begin{samepage}
\subsection{Conclusion physique}
Au froid, l'énergie rejoint le minimum exact
$u_0=-3.813$ eV/atome, $q\simeq1$ et les deux paramètres de Cowley
prennent les valeurs de $L1_2$. En chauffant, les défauts chimiques
augmentent et l'ordre global chute dans la zone du maximum de capacité
calorifique, situé à $\Tpeak$ K sur cette grille. Au chaud, $q$ devient
faible alors que les corrélations locales restent mesurables : perdre
l'ordre à longue portée ne rend pas immédiatement les occupations
indépendantes. La valeur $\Tpeak$ K caractérise donc ce modèle de réseau
rigide dans la boîte de $\Natoms$ sites ; elle n'est ni une extrapolation
à taille infinie ni une mesure expérimentale de Cu$_3$Au.
Les grandes tailles situent elles aussi leur maximum dans la zone
1050--1055 K, tout en révélant un ralentissement important. Leur ajout
consolide les tendances qualitatives ; il ne supprime pas l'incertitude
sur les poids des populations ni l'erreur de modèle.
\end{samepage}

\section{Bonus -- Histoire thermique, visualisation et corrélations spatiales}
\label{sec:bonus-cycle}
Le bonus étudie la mémoire du trajet thermique à budget de simulation
fixé. Il complète l'étude canonique précédente par un protocole
séquentiel de chauffage puis de refroidissement.
\subsection{Partie 1 -- Cycle thermique et visualisation des configurations}
\begingroup\sloppy
\IfFileExists{bonus_cycle.tex}{\input{bonus_cycle.tex}}{
Le module \texttt{run\_thermal\_cycle.py} produit les mesures et les
comparaisons des deux branches ; aucun résultat de cycle n'est joint
à cette version du rapport.}
\IfFileExists{bonus_animations.tex}{\input{bonus_animations.tex}}{}
\subsection{Partie 2 -- Longueur de corrélation et analyse de Fourier}
\IfFileExists{bonus_correlation.tex}{\input{bonus_correlation.tex}}{
Le module de corrélations produit les champs d'ordre locaux, les spectres
et les diagnostics de la longueur de second moment.}
\subsection{Partie 3 -- Analyse croisée et superpositions L6/L9/L12}
\label{sec:tailles}
\IfFileExists{bonus_sizes.tex}{\input{bonus_sizes.tex}}{
Les résultats de taille supplémentaire ne sont pas disponibles dans cette version.}
\IfFileExists{bonus_size_cycles.tex}{\input{bonus_size_cycles.tex}}{}
\endgroup

\section{Conclusion générale et perspectives}
\label{sec:conclusion-generale}
Nous avons développé et vérifié un programme répondant aux neuf questions
du sujet : géométrie CFC périodique,
échanges conservant la composition, Metropolis, fonctions \texttt{cowley}
et \texttt{ordre}, vérifications, énergie fondamentale et grandeurs
thermiques. Les bonus distinguent la mémoire du parcours, les images
de relaxation, les corrélations spatiales et les effets de taille.
Le fondamental $u_0=-3.813$ eV/atome est démontré pour le Hamiltonien,
tandis que les trois tailles situent la transformation observée vers
1050--1055 K. Au chaud, le résidu de $q$ diminue avec le volume et des
préférences chimiques locales subsistent. Ces résultats sont cohérents
entre eux, sans être tous statistiquement indépendants.

\textbf{Ce qui n'est pas encore établi.} Le maximum L6 à 1050 K est
le repère le mieux contrôlé près de la transformation, pas une température
extrapolée à volume infini. Les pics plus élevés des grandes boîtes et
la mémoire des cycles sont compatibles avec une transition du premier
ordre, mais les poids des populations critiques restent insuffisamment
explorés, particulièrement en L12. Ni une courbe arrondie ni un creux
de Binder ne tranchent seuls entre premier et second ordre. Le caractère
du premier ordre du matériau réel et le saut du champ moyen ne prouvent
pas celui de notre Hamiltonien spatial. Enfin, les grandes longueurs
spectrales froides ne sont pas des rayons de domaines, et les GIF ne
mesurent pas une cinétique physique de nucléation.

\subsection*{Priorité 1 : améliorer l'exploration avant d'agrandir encore}
Il faut d'abord prolonger les productions critiques L9/L12 et ajouter
des graines indépendantes, avec des départs ordonnés et désordonnés.
On recherchera des passages répétés entre les populations, des proportions
de séjour reproductibles et une stabilité des moyennes lorsque la durée
augmente. ESS, R-hat, blocs, demi-séries et circulation des répliques
doivent être examinés ensemble ; un ESS additionné élevé peut masquer
une chaîne lente. La grille sera adaptée aux liens limitants et à la
circulation globale. Un meilleur taux d'échange seul n'est pas le critère
d'arrêt. Si les barrières restent limitantes, un échantillonnage étendu
avec poids corrigés et validation sur une petite boîte pourra être étudié.

\subsection*{Priorité 2 : tester la nature de la transition à précision comparable}
Une étude de taille doit viser une qualité d'exploration comparable,
plutôt qu'un nombre identique de balayages. À proximité de la coexistence,
on examinera les histogrammes joints d'énergie par atome et d'ordre,
séparément par chaîne, puis leur stabilité. Pour soutenir un premier
ordre, on cherchera notamment une séparation énergétique entre populations
qui reste non nulle lorsque $L$ augmente : elle donnerait une chaleur
latente dans la limite thermodynamique. Dans le régime asymptotique et
à coexistence correctement échantillonnée, un tel saut entraîne une
croissance du maximum de $c_V$ \emph{par atome} proportionnelle à $N$.
Une transition continue demande d'autres lois de taille et un contrôle
des corrélations critiques ; un pic de chaleur spécifique n'y est pas
nécessairement divergent. Historiques, histogrammes, cumulants et
lois de taille doivent donc être croisés \cite{challa}.
On n'applique pas ces critères quantitativement aux trois campagnes
actuelles : leurs budgets et leurs diagnostics diffèrent, et quelques
points ne suffisent pas à établir le régime asymptotique.

\subsection*{Priorité 3 : séparer mémoire thermique, domaines et cinétique}
Les cycles L6/L9/L12 sont déjà terminés. La suite utile consiste à
les répéter avec davantage de graines et plusieurs durées de palier,
en conservant l'appariement chauffage--refroidissement de chaque cycle.
La dépendance des écarts à la durée permettra de tester la relaxation
insuffisante ; deux répétitions ne donnent pas une largeur d'hystérésis
robuste. Pour les corrélations, on pourra comparer les directions et
attribuer une incertitude au diagnostic intra-variante. Pour les GIF,
il faudra identifier des domaines locaux et leurs interfaces, répéter
les trempes et conserver une résolution temporelle adaptée. Une étude
de nucléation physique demanderait en plus une dynamique locale,
par exemple atome--lacune, des barrières de migration et une calibration
du temps : les échanges globaux actuels restent un outil d'échantillonnage.

\subsection*{Priorité 4 : rapprocher le modèle du matériau}
Le décalage avec Cu$_3$Au réel doit être étudié en modifiant séparément
les hypothèses physiques : paramètres d'interaction validés,
relaxations atomiques, vibrations et variation de volume.
Une comparaison calorimétrique demanderait de distinguer la contribution
configuratoire calculée ici du signal expérimental total et de son
protocole. Augmenter seulement $L$ ne rétablit aucune de ces contributions.
Le projet établit ainsi un modèle explicatif de l'ordre chimique ; les
prolongements prioritaires visent d'abord la fiabilité statistique,
puis la caractérisation thermodynamique et la fidélité au matériau.

\appendix
\section{Correspondance entre les questions et les fichiers}
\begin{center}\small
\begin{tabular}{p{.12\linewidth}p{.40\linewidth}p{.38\linewidth}}
\toprule
Question & Réponse numérique & Fichier principal\\
\midrule
1 & Réseau, énergie, échanges, Metropolis & \texttt{lattice.py}, \texttt{model.py}, \texttt{kernel.py}\\
2 & Cowley des deux coquilles & \texttt{model.py : cowley}\\
3 & Ordre signé, amplitude, sous-réseaux & \texttt{model.py : ordre}\\
4 & Tests analytiques et probabilistes & \texttt{tests/test\_physics.py}\\
5 & Énergie exacte, preuve et tests & \texttt{model.py}, présent rapport\\
6 & Configurations, répliques, graines & \texttt{simulation.py}, \texttt{chain\_*.npz}\\
7 & Moyennes, erreurs, figures & \texttt{statistics.py}, \texttt{figures.py}\\
8 & Maximum sur la grille, champ moyen & \texttt{summary.json}, \texttt{meanfield.py}\\
9 & Discussion fondée sur les données & Rapport et diagnostics CSV\\
10 & Cycle chauffage--refroidissement & \texttt{thermal\_cycle.py}, \texttt{run\_thermal\_cycle.py}\\
\bottomrule
\end{tabular}\end{center}

Le lanceur \texttt{run\_project.py} enchaîne les tests, la simulation, les
statistiques, les figures et le rapport. L'option \texttt{--reuse} réutilise
uniquement une campagne dont la configuration correspond exactement.
La compilation PDF demande une installation LaTeX, mais les fichiers
\texttt{.tex} et les figures sont produits même sans compilation.

\begin{thebibliography}{9}
\bibitem{metropolis} N. Metropolis, A. W. Rosenbluth, M. N. Rosenbluth,
A. H. Teller et E. Teller, \emph{Equation of State Calculations by Fast
Computing Machines}, J. Chem. Phys. \textbf{21}, 1087--1092 (1953).
\href{https://doi.org/10.1063/1.1699114}{doi:10.1063/1.1699114}.
\bibitem{hukushima} K. Hukushima et K. Nemoto, \emph{Exchange Monte Carlo
Method and Application to Spin Glass Simulations}, J. Phys. Soc. Jpn.
\textbf{65}, 1604--1608 (1996).
\href{https://doi.org/10.1143/JPSJ.65.1604}{doi:10.1143/JPSJ.65.1604}.
\bibitem{binder} K. Binder, \emph{Theory of first-order phase transitions},
Rep. Prog. Phys. \textbf{50}, 783 (1987).
\href{https://doi.org/10.1088/0034-4885/50/7/001}{doi:10.1088/0034-4885/50/7/001}.
\bibitem{challa} M. S. S. Challa, D. P. Landau et K. Binder,
\emph{Finite-size effects at temperature-driven first-order transitions},
Phys. Rev. B \textbf{34}, 1841--1852 (1986).
\href{https://doi.org/10.1103/PhysRevB.34.1841}{doi:10.1103/PhysRevB.34.1841}.
\bibitem{benisek} A. Benisek et E. Dachs, \emph{The vibrational and configurational
entropy of disordering in Cu$_3$Au}, Journal of Alloys and Compounds
\textbf{632}, 585--590 (2015).
\href{https://doi.org/10.1016/j.jallcom.2014.12.215}{doi:10.1016/j.jallcom.2014.12.215}.
\bibitem{owen} L. R. Owen, H. Y. Playford, H. J. Stone et M. G. Tucker,
\emph{Analysis of short-range order in Cu$_3$Au using X-ray pair distribution functions}.
\href{https://doi.org/10.17863/CAM.38392}{Manuscrit accepté, dépôt de Cambridge}.
\bibitem{kessler} M. Kessler, W. Dieterich et A. Majhofer,
\emph{Ordering kinetics in an fcc A$_3$B binary alloy model: Monte Carlo studies},
Phys. Rev. B \textbf{67}, 134201 (2003).
\href{https://arxiv.org/abs/cond-mat/0207611}{Prépublication et protocole de sauts atome--lacune}.
\bibitem{paras} J. Paras et A. Allanore, \emph{Contribution of electronic
entropy to the order-disorder transition of Cu$_3$Au}, Phys. Rev. Research
\textbf{3}, 023239 (2021).
\href{https://doi.org/10.1103/PhysRevResearch.3.023239}{doi:10.1103/PhysRevResearch.3.023239}.
\end{thebibliography}
\end{document}
